% from ega4-2-fr.tex
\oldpage[IV-2]{212}

de $A$ soit réduit et que, si $\mathfrak q_j$ ($1\leq j\leq n$) sont les
idéaux premiers minimaux de $\widehat A$, les corps des fractions des anneaux
intègres $\widehat A/\mathfrak q_j$ soient des extensions séparables de $K$.
\end{corollary}

\begin{proof}
Les quotients intègres de $A$ sont en effet $A$ lui-même et des corps, donc la
condition $b)$ de (7.6.4) équivaut à l'hypothèse que $A$ est un anneau
japonais, et l'hypothèse $a')$ aux conditions de l'énoncé sur $\widehat A$,
d'où la conclusion.
\end{proof}

\begin{remarks}[7.6.7]
\label{IV.2.s7.6.7-fr}
\begin{enumerate}
\item[\emph{(i)}] Les conditions équivalentes de (7.6.6) signifient encore que
les fibres formelles de $A$ sont géométriquement régulières (7.3.19, (iv)). On
a déjà observé (\emph{loc. cit.}) que cette condition est vérifiée lorsque $A$
est un anneau de valuation discrète dont le corps des fractions est de
caractéristique 0.

\item[\emph{(ii)}] Le même raisonnement que dans la première partie de la
démonstration de (7.6.4) montre que pour un anneau semi-local noethérien $A$,
les deux conditions suivantes sont équivalentes :
\begin{enumerate}
\item[\emph{a)}] Pour tout anneau quotient intègre $B$ de $A$, le complété
$\widehat B$ est réduit.
\item[\emph{a$'$)}] Les fibres formelles de $A$ sont réduites.
\end{enumerate}
Compte tenu de (0, 23.1.7, (i)), ces deux conditions impliquent la suivante :
\begin{enumerate}
\item[\emph{b)}] Pour tout anneau quotient intègre $B$ de $A$, la clôture
intégrale de $B$ est une $B$-algèbre finie.
\end{enumerate}
Lorsque $A$ est un anneau universellement caténaire, on peut prouver que la
condition $b)$ entraîne inversement $a)$; nous n'aurons pas à utiliser ce
résultat.
\end{enumerate}
\end{remarks}

\subsection{Applications : II. Anneaux universellement japonais}
\label{IV.2.ss7.7-fr}

\begin{env}[7.7.1]
\label{IV.2.s7.7.1-fr}
Rappelons (0, 23.1.1) que l'on appelle anneau universellement japonais un
anneau $A$ tel que toute $A$-algèbre intègre de type fini soit un anneau
japonais. Il revient au même de dire que pour toute $A$-algèbre intègre de type
fini $B$, la clôture intégrale de $B$ est une $B$-algèbre finie.
\end{env}

\begin{theorem}[7.7.2]
\EGAReaderTarget{ega4.c04.thm.7-7-2}{IV.2.s7.7.2-fr}
(Nagata). --- Soit $A$ un anneau noethérien. Les conditions suivantes sont
équivalentes :
\begin{enumerate}
\item[\emph{a)}] $A$ est un anneau universellement japonais.
\item[\emph{b)}] Tout anneau quotient intègre de $A$ est un anneau japonais.
\item[\emph{c)}] Pour tout idéal maximal $\mathfrak m$ de $A$, tout anneau
quotient intègre de $A_{\mathfrak m}$ est un anneau japonais, et pour tout
anneau quotient intègre $B$ de $A$, de corps des fractions $K$, toute
extension radicielle finie $K'$ de $K$ et toute sous-$B$-algèbre finie $B'$ de
$K'$, de corps des fractions $K'$, il existe $f'\neq0$ dans $B'$ tel que
$B'_{f'}$ soit intégralement clos (cf. (6.13.7)).
\end{enumerate}
De plus, si $A$ vérifie ces conditions, il en est de même de tout anneau de
fractions $S^{-1}A$ de $A$ et de toute $A$-algèbre de type fini.
\end{theorem}

\begin{proof}
Montrons d'abord que $b)$ entraîne $c)$; tout anneau quotient intègre de
$A_{\mathfrak m}$ est un anneau de fractions d'un anneau quotient intègre de
$A$, donc un anneau japonais (0, 23.1.1). D'autre part, tout quotient $B$ de
$A$ par un de ses idéaux premiers est un anneau japonais, donc il en est de
même de $B'$ (0, 23.1.1). On en déduit que l'ensemble
$\operatorname{Nor}(\operatorname{Spec}(B'))$ est ouvert (6.13.3), donc
contient un ensemble ouvert non vide
$D(f')=\operatorname{Spec}(B'_{f'})$.

